\textbf{Choice of method.} Background augmentation accounted for most of the observed reduction in movement flips. A paired reconstruction target directly implements our hypothesis about relevance. It needs aligned views that retain the handwriting while changing the decoration. The Gaussian needs only a center and a width. Across the test set its mean normalized pixel weight was 1.79 on ink and 0.67 on the mark. Even in the outermost quartile of digit centroids, its mean marked ACC was 73.11\% compared with 71.14\% under augmentation alone. The farthest MNIST test centroid was only 1.41 pixels from the image center. This subgroup cannot establish robustness to substantially displaced digits.

\textbf{Limits of the relevance assumption.} In this construction all stroke pixels stay fixed and the two mark sites are disjoint, so equality across the paired views yields a clean target. Repeated background features could also be stable, and pose changes could move relevant foreground pixels. A fixed Gaussian can downweight an off-center stroke or upweight a central distractor. Site selection here avoids the handwriting, so where a mark can appear depends on the available blank space. The findings cover this family of decorations, this split and five seeds. A separate test with displaced digits and backgrounds near the image center would probe the limits of the spatial prior.

\textbf{Relation to the earlier experiments.} The \href{https://github.com/JohnR31N/cs126FinalProj/blob/idec-final-v1.4/IDEC/report/report.pdf}{additive-overlay study} tested a known, fixed stripe direction and crossed reconstruction losses across the two training stages. Those stripes changed intensities on the digit as well as in the surrounding space. Here the foreground pixels stay exactly fixed, and the paired target is constructed without the shape masks or a known stripe direction. Each study has its own matched baseline. The difference in their numerical effect sizes also reflects the extent, strength and type of intervention.
